\documentclass[11pt,a4paper]{article}
\usepackage[utf8]{inputenc}
\usepackage[T1]{fontenc}
\usepackage[spanish]{babel}
\usepackage{hyperref}
\usepackage{geometry}
\geometry{margin=2.5cm}
\usepackage{parskip}

\title{090 la cita se puede usar, la marca no}
\author{Alejandro Valenzuela Rivera\\\small Universidad Autónoma de Nuevo León}
\date{\today}
\begin{document}
\maketitle

\begin{quote}
\textbf{Sesgo declarado.} Este trabajo se sostiene desde el sesgo declarado de su autor, y no finge neutralidad. Lo que afirma se puede tumbar, y ese es el punto: el falsificador está escrito, no escondido.
\end{quote}

\section*{Resumen}
Citar una obra para pensar algo es legitimo y ni siquiera necesita permiso. Usar el nombre de la obra o el nombre del estudio como nombre del sistema es otra cosa: es tomar una marca, y las marcas no se citan, se licencian.

\section{Por qué se sostiene}
La linea es esta y es clara: la referencia se declara una vez en la pieza, con su obra y sus autores; el nombre propio del sistema no puede ser el nombre de la obra ni el de su estudio. Que un nombre sea una referencia evidente no es defensa: es lo que hace que se note. Y esto importa el triple en un repo que es de la gente y que va a leerse, porque una marca ajena en el titulo convierte la obra en un problema legal para todos los que la usen.

\section{De dónde viene}
Gurren Lagann (Gainax; derechos en Aniplex) y Studio Trigger.

\section{El caso}
Gurren Lagann y Studio Trigger.

\section{Método}
Una cita de la obra para analizarla, con su autor nombrado.

\section{Límites, y qué lo tumba}
El nombre de una obra usado como nombre propio del sistema.

\section{Lo que queda abierto}
\textquestiondown{}Si me inspiro en una obra, mi sistema puede llamarse como ella?

\section*{Cómo se leyó esta tesis}
Esto se escribió en el repositorio man-made-empathy, del que este candado es una pieza: candados/03 el nombre y la obra/090 la cita se puede usar, la marca no.md. La máquina que lo compuso lo hizo desde este mismo archivo, y no agregó nada que no estuviera escrito. Si algo de esto es falso, se dice en el mismo sitio, y se corrige con fecha.

\end{document}
